\section{The orientation generator and the path-fibration proof}
\label{mapsI:course:homotopy-section}

The homotopy-sphere conclusion is classical; the following argument retains
our original map and gives the distinct path-fibration proof of the course.
The complete elementary constructions used below are printed in
\cref{mapsI:cw:section,mapsI:pf:section}.  The first Hurewicz theorem and
homological Whitehead theorem are also stated in
\cite[Theorem~4.32 and Corollary~4.33]{Hatcher2002}; we use their integral,
simply connected forms.

\begin{theorem}[the original degree-one homotopy equivalence]
\label{mapsI:course:CG-S6-07-HOMOTOPY}
For the constructed threefold, with its complex orientation, there is a
map \(g:S^6\to X\) satisfying \(g_*[S^6]=[X]\), and this particular map
is a homotopy equivalence.
\end{theorem}

\begin{proof}
The fundamental-group calculation in
\cref{mapsI:top:prop:geometric-pi1} gives \(\pi_1(X)=0\), and
\cref{mapsI:top:thm:h2-vanishing,mapsI:top:lem:cohomological-closure}
give the complete integral homology of the sphere.  For \(2\le n\le5\),
assume inductively that \(\pi_i(X)=0\) for \(1\le i<n\).  The integral
Hurewicz map
\[
 h_n:\pi_n(X)\longrightarrow H_n(X;\mathbb Z),
 \qquad [a]\longmapsto a_*[S^n],
\]
is then an isomorphism, by \cref{mapsI:cw:hurewicz-theorem-result}.
The target vanishes, so the induction gives \(\pi_i(X)=0\) for
\(1\le i\le5\).  In degree six choose \([g]\) mapping to the positive
orientation generator \([X]\).  Thus
\begin{equation}
 g_*[S^6]=[X].
 \label{mapsI:course:degree-one}
\end{equation}
This map induces an integral homology isomorphism in every degree:
in degrees zero and six it preserves the generators, and all remaining
groups on both sides vanish.

Replace \(g\) by its actual path fibration
\[
 E_g=\{(s,\lambda):s\in S^6,\ \lambda(0)=g(s)\},
 \qquad p(s,\lambda)=\lambda(1).
\]
The path space has the compact-open topology.  The homotopy
\((s,\lambda)\mapsto(s,u\mapsto\lambda((1-t)u))\) strongly retracts
\(E_g\) onto the constant paths, identified with \(S^6\).
Under this retraction \(p\) represents \(g\).
For clarity, the lifting formula used in \cref{mapsI:pf:path-lifting}
is the following.  Given \(H:Z\times I\to X\) and an initial lift
\((s(z),\lambda_z)\), the path at time \(t\) is
\[
 \eta_{z,t}(u)=
 \begin{cases}
 \lambda_z((1+t)u),&0\le u\le (1+t)^{-1},\\
 H(z,(1+t)u-1),&(1+t)^{-1}\le u\le1.
 \end{cases}
\]
The pieces agree at the join, their two closed domains cover the parameter
space, and compact-open adjunction proves continuity.  The last endpoint
is \(H(z,t)\), while the time-zero path is exactly \(\lambda_z\).
This proves the homotopy lifting property.  Relative cell lifting gives
the exact homotopy sequence, as proved in
\cref{mapsI:pf:path-lifting}.

For the cellular filtration use the actual comparison constructed in
\cref{mapsI:cw:cw-manifolds}:
\[
 \alpha:T_X\longrightarrow X,\qquad
 \beta:X\longrightarrow T_X,\qquad
 \alpha\beta=1_X,\qquad \beta\alpha\simeq1_{T_X}.
\]
Put \(\bar g=\beta g\), let \(\bar E=E_{\bar g}\), and write
\(\bar p(s,\lambda)=\lambda(1)\).  There is an explicit comparison
\[
 J:\bar E\longrightarrow E_g,\qquad
 J(s,\lambda)=(s,\alpha\circ\lambda),
 \qquad pJ=\alpha\bar p.
\]
Its starting point is correct because
\(\alpha\beta g(s)=g(s)\), exactly.  Both total spaces retract onto the
same \(S^6\), and \(J\) commutes with these retractions.  The base
\(T_X\) is an actual simply connected CW complex with \(\pi_2(T_X)=0\).
The map \(\bar p\) is an integral homology isomorphism in every degree,
since \(\bar g=\beta g\) is one.  Let \(K\) be its fibre over a vertex.
It is path connected: a path between two fibre points in the connected
total space projects to a null-homotopic loop, and lifting its contraction
relative to the two endpoints changes that path to one in the fibre.
The exact sequence, together with \(\pi_2(T_X)=0\) and
\(\pi_1(\bar E)=\pi_1(S^6)=0\), then gives \(\pi_1(K)=0\).

Suppose a positive homotopy group of \(K\) is nonzero and choose the least
such degree \(n\ge2\).  Hurewicz gives
\[
 H_i(K;\mathbb Z)=0\quad(0<i<n),\qquad
 H_n(K;\mathbb Z)\cong\pi_n(K)\ne0.
\]
The integral homology spectral sequence for \(\bar p\), constructed
with its actual edge maps in \cref{mapsI:pf:filtered-chains,mapsI:pf:serre-homology},
has constant coefficient systems because the base is simply connected.
In total degree \(n\), the only possibly nonzero terms are
\[
 E_2^{0,n}=H_n(K),\qquad E_2^{n,0}=H_n(T_X).
\]
No differential leaves column zero.  A differential entering \((0,n)\)
has source \((r,n-r+1)\).  The sources for \(2\le r\le n\) vanish by
the displayed fibre-homology vanishings, and for \(r>n+1\) they have
negative second degree.  The sole possible incoming map is therefore
\[
 d_{n+1}:E_{n+1}^{n+1,0}\longrightarrow E_{n+1}^{0,n}.
\]
The source is still \(H_{n+1}(T_X)\): there are no incoming maps to
the bottom row, and all earlier outgoing targets have fibre degree
between one and \(n-1\).  The actual edge image
\(\operatorname{im}(\bar p_*:H_{n+1}(\bar E)\to H_{n+1}(T_X))\)
is \(E_\infty^{n+1,0}=\ker d_{n+1}\).  Surjectivity of \(\bar p_*\)
makes this kernel the whole source, so \(d_{n+1}=0\).
Consequently
\[
 H_n(K)=E_\infty^{0,n}=F_0H_n(\bar E).
\]
The other actual edge map \(\bar p_*:H_n(\bar E)\to H_n(T_X)\)
kills this first filtration subgroup: its kernel is
\(F_{n-1}H_n(\bar E)\), which contains \(F_0\) since \(n\ge2\).
Injectivity of this edge map therefore forces \(H_n(K)=0\), a
contradiction.  All homotopy groups of \(K\) vanish.

The exact homotopy sequence makes \(\bar g\), and hence the original
\(g\), an isomorphism on every homotopy group.  The complete CW-domain
compression argument in \cref{mapsI:pf:cw-inverses} supplies an actual
inverse \(\bar h:T_X\to S^6\) to \(\bar g\), with both inverse
homotopies.  Set
\[
 h_X=\bar h\beta:X\longrightarrow S^6.
\]
Then \(h_Xg=\bar h\bar g\simeq1_{S^6}\).  For the other composition,
use \(g=\alpha\bar g\), then \(\bar g\bar h\simeq1_{T_X}\), then
the exact equality \(\alpha\beta=1_X\), to obtain
\[
 gh_X=\alpha\bar g\bar h\beta\simeq\alpha\beta=1_X.
\]
Both homotopies concern the original spaces and the original degree-one
map.  The orientation generator in \eqref{mapsI:course:degree-one}
has not been rescaled.
\end{proof}
